\section{Introduction}

Autonomous multi-agent systems are increasingly deployed in roles that carry genuine consequential weight: allocating resources, composing and deploying code, initiating communications, and making downstream decisions that affect human welfare. As the operational scope of these systems expands, so too does the urgency of a question that the field has largely deferred: \emph{what happens when things go wrong, and who is empowered to stop it?} We argue that the answer cannot be a policy statement, a warning label, or a human-in-the-loop checkpoint that is optional by construction. Accountability in autonomous multi-agent systems must be encoded as a structural fallback, a hardwired termination pathway that activates when the system's behavior diverges from its \bxthree\ in ways that other mechanisms cannot or will not correct.

The problem is not hypothetical. well-documented failures in production agentic systems reveal a consistent pattern: cascades of action propagate faster than human oversight can intervene, error states persist because each agent in a chain is individually within its local operating parameters, and the absence of a system-wide termination trigger allows consequential actions to compound before anyone possesses the authority or the mechanism to halt them \cite[]{}. In工业 and infrastructural deployments, the consequences range from data corruption and resource exhaustion to physical process deviations that require manual emergency shutdown. The root cause is architectural: most deployed multi-agent systems treat bail-out as an afterthought, a capability to be considered only after the primary operational logic is complete. This ordering is backwards. When a system's potential for harm scales with its autonomy, the existence of a mandatory, architecturally enforced termination pathway is not a feature among features—it is a pre-condition for responsible deployment.

The \bxthree\ framework (\textbf{B}ounds, \textbf{X}change, \textbf{H}eight, \textbf{3}) provides a conceptual foundation for reasoning about this problem. Every agentic system is oriented by a \purpose: an operational objective that defines what the system is attempting to accomplish. \Bounds\ define the envelope of permissible action—the ethical, legal, and operational constraints that mark the difference between acceptable and unacceptable behavior. \Fact\ is the verification standard: claims made by the system or about the system must be traceable to actionable evidence, not inferred or approximated. In this framing, a bail-out mechanism is not an additional feature bolted onto an existing architecture; it is the architectural embodiment of \bounds, the physical proof that constraints are enforceable rather than advisory.

Consider the alternative. A system whose \bounds\ are encoded only as policy—language in a terms-of-service document, a clause in a system prompt—has no mechanism to enforce those bounds when the system itself is the agent of action. When a multi-agent cascade is underway, and each successive agent is acting within its individual \purpose, the aggregate behavior may violate system-level \bounds\ without any single agent perceiving the violation. Policy-encoded constraints fail precisely at the moment when they are most needed: during autonomous operation, under load, when the system's outputs are propagating to other systems or to the physical world. Accountability without an architectural fallback is accountability in name only.

The core insight of this paper is that \textbf{accountability must have an architectural fallback}. By this we mean that the termination of a multi-agent system—when it exceeds its \bounds, when its actions can no longer be reconciled with \fact, or when its \purpose\ has become unachievable or harmful—must be executable by the system itself, not solely by external human intervention. The fallback must be \mandatory: once the triggering conditions are met, termination proceeds regardless of whether any human is monitoring. The fallback must be \hierarchical: it overrides all other operational directives, including any agent's local \purpose, because \bounds\ at the system level supersede component-level objectives. And the fallback must be \verifiable: its activation is logged against \fact, creating an auditable record of why the system was terminated and under what conditions. These three properties—mandatory enforcement, hierarchical priority, and verifiability—distinguish a genuine bail-out mechanism from a graceful shutdown option or a kill switch that can be ignored.

This paper makes the following contributions. We first establish the theoretical necessity of mandatory bail-out in autonomous multi-agent systems by analyzing failure modes that arise specifically from the absence of such a mechanism (\S\ref{sec:problem}). We then present the \bxthree\ framework as a principled basis for defining bail-out conditions in terms of \purpose, \bounds, and \fact, distinguishing our approach from shutdown switches and circuit breakers that lack hierarchical priority (\S\ref{sec:framework}). We describe the Bailout Protocol itself: its triggering conditions, its enforcement hierarchy, its fallback behavior, and its logging semantics (\S\ref{sec:protocol}). We evaluate the protocol against a suite of adversarial and accidental failure scenarios, demonstrating that it terminates dangerous behavior that existing approaches permit to persist (\S\ref{sec:evaluation}). Finally, we discuss related work in agent safety, multi-agent coordination, and accountability architecture, and identify open problems for future research (\S\ref{sec:related}).

The remainder of this paper is organized as follows. Section~\ref{sec:problem} examines the failure modes that motivate the Bailout Protocol. Section~\ref{sec:framework} introduces the \bxthree\ terminology and its application to bail-out condition specification. Section~\ref{sec:protocol} details the protocol specification. Section~\ref{sec:evaluation} presents the evaluation. Section~\ref{sec:related} discusses related work and open questions.

\label{sec:introduction}
